可观测性把运行中的症状连接到持久事实与恢复负责人。一份付款可以在 HTTP trace 结束数小时后继续处理，业务记录负责长期关联，trace、日志与指标解释某次执行。运维接口沿这些关联读取证据，再通过受权威核验的命令推进责任。

\subsection{从 S 走到各保存域的事实}

完整 $S$ 是业务入口，$A_J$ 保存原输入、P/C/T 和原平台范围与键。组合购买从 C 关联 TX，再查询参与者的 prepared、票据和决定应用状态。平台成功对象 X 与 $R_J$ 关联，原订单 O 关联 EV，消费者通过 EV 与业务守卫找到 F，F 再关联 dispatch 和外部发运对象。RF 独立关联原收款和退款占额。图~\ref{fig:ch21-observe} 的连线表示这些持久关联。trace 关联一次请求经过的局部调用，attempt 标记其中某次执行尝试，它们帮助定位耗时与错误，跨重试的业务关联则从原 S、C 和 TX 读取。

\CTwentyOneFigure{ch21-09}{持久身份连接跨域业务事实，trace 和 attempt 描述其中一次执行}{fig:ch21-observe}

观察服务按字段聚合 G1 的结果、平台收款和消费库进度，每个字段注明来源及读取时点。这些查询各自在所属保存域执行，返回时点可以不同。平台先完成、权威后回填的几秒钟内，视图可能同时显示“平台成功”和“权威未知”。界面把两份观察并列展示，恢复器便能沿平台对象定位尚未完成的权威回填。

下面这份观察对象由身份解析服务组装。服务先按调用方权限找到原记录，再从权威、库存、平台和调度器读取观察，映射到相应字段，并保存各自来源与读取时间：
\begin{BookCode}
{
  "identity": {"account_id": 1,
    "action": "purchase.create.external.v2",
    "session_id": "E1", "idempotency_key": "K1"},
  "authority": {"partition": "b1", "owner_epoch": 1,
    "purchase_id": "P:b1:1",
    "charge_id": "C:b1:1", "task_id": "T:b1:1",
    "binding": "accepted", "result": null,
    "source": "G1_query", "observed_at": "2026-10-18T01:00:02Z"},
  "inventory": {"tx_id": "TX:P:b1:1", "decision": "commit",
    "participants_applied": ["b2", "b3"],
    "source": "TX_query", "observed_at": "2026-10-18T01:00:03Z"},
  "payment": {"scope": "simulator:merchant-7:test:charge.v1",
    "key": "pay-C:b1:1", "observed_status": "succeeded",
    "provider_id": "X1", "amount_minor": 3000,
    "currency": "CNY", "source": "provider_query",
    "observed_at": "2026-10-18T01:00:04Z"},
  "recovery": {"phase": "awaiting_authority_observation",
    "worker_generation": 8,
    "next_action": "validate_and_record_observation",
    "source": "scheduler", "observed_at": "2026-10-18T01:00:05Z"}
}
\end{BookCode}

这份记录已经定位到一个具体缺口：库存决定已提交且参与者已应用，平台查询取得 X1，购买权威尚未保存原结果。恢复者核对 X1 的原账户、金额、币种及 C 关联，再用当前代次向权威提交观察。如果权威已由回调先一步完成，同样的查询观察返回已保存结果；若编号或金额不同，接口返回冲突并保留两份证据。聚合服务负责读取，结果写入经过权威的资格、原输入和已有结果检查。

\subsection{原结果和当前进度分开读取}

\code{result} 读取不可覆盖的原业务响应 $R(S)$，其中保存原受理内容、付款结果及实际订单编号。实现中这一响应的完整字段为 \code{acceptance}、\code{payment\_result}、保存时的 \code{progress} 和 \mbox{\code{result\_version}}。这里保存的进度描述原响应形成时的事实。$R_J$ 另外固定完整 \code{provider\_result}、付款结果、O 与 EV，受授权的证据查询沿原 C 读取这些对象。原业务响应回答这份购买当时形成了什么确定结果，完整平台对象从证据接口取得。

当前 \code{progress} 观察读取最新责任阶段、最近一次可信观察、恢复计划和履约、退款等后续关联，数据可以继续变化。原购买成功以后发生退款，原结果继续说明原收款成功，退款进度从关联 RF 读取。把原响应与当前观察分别提供，客户端既能重放当时的确定结果，也能展示现在的业务进展。

例如 $R(S)$ 已保存 X1、3000 分、O1 和付款成功。后来履约平台不可达，dispatch 未知，界面可以同时展示“原购买已付款”和“发货核对中”。再后来退款 RF1 成功 1000 分，原购买回放仍是原响应，当前页面另外展示已退金额和剩余额度。把 $R(S)$ 改写成“退款中”会让同一购买重试得到不同含义，也丢掉原事实与后续动作之间的关系。

进度响应携带最近观察时间、权威应用位置或投影版本，客户端据此决定继续展示还是转向权威。日志中的候选 P、暂时生成的对象 ID 和最终已提交编号分别标注。例如事务尝试先打印候选 P99，随后回滚，重试命中原绑定 P17。观察接口沿权威返回的 P17 关联 C、O 与 EV，排错时再用 P99 查找那次回滚尝试。

权限也按事实分层。普通用户查看自己账户的原结果及经脱敏的进度，客服查看授权业务范围内的订单、平台对象和恢复原因，值班工程师查看组内日志与调度状态，财务对账角色提交平台证据。平台密钥、完整个人信息和付款凭据按其保存规则读取，业务日志只放受控的定位摘要。摘要帮助索引，完整输入比较从受保护的原绑定读取。

\subsection{日志怎样保留因果与拒绝原因}

结构化日志为每次局部动作记录时间、服务、分区、owner epoch、Raft 任期、工作代次、trace、attempt 及当前已知的持久身份。尚未形成 O、EV 或 F 时填空值，后续事件补充关联。平台查询返回以后，观察要经过权威提交，响应再经过网络送达客户端，三个事件分别记录。响应丢失时，前两个事件已经存在，客户端从原 S 查询才能获知结果。下面抽出平台观察、权威提交和旧代次拒绝这三类日志的业务关联字段。

\begin{BookCode}
{"event":"payment_query_observed","charge_id":"C:b1:1",
 "attempt":"pay-attempt-19","worker_generation":8,
 "scope":"simulator:merchant-7:test:charge.v1","key":"pay-C:b1:1",
 "provider_id":"X1","provider_status":"succeeded",
 "amount_minor":3000,"currency":"CNY"}
{"event":"payment_observation_committed","charge_id":"C:b1:1",
 "worker_generation":8,"provider_id":"X1",
 "order_id":"O:b1:1","event_id":"EV:b1:1"}
{"event":"payment_observation_rejected","charge_id":"C:b1:1",
 "worker_generation":7,"current_generation":8,
 "provider_id":"X1","reason":"stale_worker_generation"}
\end{BookCode}

这三条日志可以来自不同进程和不同 trace。代次 7 的拒绝说明旧执行者失去写入资格，代次 8 的提交说明同一 X1 已被接受。把最后一条仅计作付款失败，会夸大失败率并误导人工重付。平台尝试状态、权威提交状态与业务最终状态分别形成指标。

跨主机时钟可能有偏差，因果关系优先取自 TX 决定、日志位置、调用父子关系和已提交关联。墙钟用于耗时和年龄估计，异常跳变时标注时钟健康。审计记录保存原始证据的接收时间、来源签名或校验结果以及记录哈希，事后可以复查一份付款观察为什么被接纳。

\subsection{指标和告警以未决责任为中心}

业务指标按新绑定、已有身份读取、异参冲突、业务拒绝、成功与未知分别计数。复制指标按组记录主副本、提交和应用位置、跟随滞后、磁盘余量、快照成功及恢复耗时。责任指标按 T、TX、EV、RF 和 dispatch 记录存量、字节、年龄分布、到期分布、流入、完成和转人工数量。入口流入下降而提交稳定时，值守者沿业务流量检查；提交变慢并伴随复制滞后时，转向组内磁盘和网络；提交稳定而 T 增长时，再沿平台与恢复调度检查。

单一成功率掩盖正在累积的未知。若一分钟受理一万份，9900 份迅速成功，100 份长期未知，成功率看似 99\%，未决资金责任却每分钟增加 100 份。监控同时显示新未知进入率、确定完成率、最老年龄与接近证据期限的数量。已转人工的 C 继续计入未终结责任，转队列动作只改变负责人。

高基数字段适合日志和查询索引，指标标签使用分区、平台范围、动作与错误类别等受控维度。把每个 $S$ 放进时序指标标签，会产生与业务量同阶的序列数，反过来压垮监控。告警附少量样例 C/TX/EV，用户从样例进入按权限查询的责任明细。

以下三类告警各自对应下一步动作。prepared 超过预算且协调组不可达时，参与者保留占用，值班恢复协调组与决定查询，取得持久决定后再应用。平台成功对象已取得而权威结果为空，恢复器检查回填池、资格和完整输入冲突。EV 已得到发布确认而消费侧迟迟无 I，值班按 broker 交付位置、消费组和消费事务检查。告警附当前事实、负责人、证据期限与样例身份，值班从样例进入相应责任的查询和恢复接口。

恢复预计完成时间也有价值。队列存量下降而最老年龄不降，可能是调度器一直跳过某类反复失败的责任；存量稳定而新进入率和完成率都很高，可能处于持续平衡。运维按 21.2 的净排空率估计恢复终点，并设置相应放量计划。调度记录保存每次领取时的年龄、证据期限和所属份额，值守者据此定位被持续跳过的类别，并调整保留预算或转入人工处理。

\subsection{控制面命令的提交与生效}

运维至少区分暂停新受理、暂停新的平台发送、允许原键查询、暂停 Outbox 发布和暂停外部履约。控制这些动作分别影响责任流入和排空。平台发生重复响应异常时，可以先暂停新的发送，保留查询、可信回填及原结果读取；消息集群故障时，可以收紧新购买，并让已确定付款继续保存在 Outbox。

\CTwentyOneFigure{ch21-10}{开关分别作用于新增责任和各类恢复动作，版本确认与在途请求一同观测}{fig:ch21-switches}

控制命令保存命令 ID、作用域、配置版本、操作者、授权角色、原因和生效条件。接口支持按期望版本更新，防止两位操作人员先后覆盖彼此的控制。例如当前 \code{merchant-7} 发送策略版本 41，命令 M9 要求在版本 41 上暂停发送并生成版本 42。另一条基于版本 41 的放开命令得到版本冲突，操作者读取 42 后重新判断。

工作者在实际发送前检查当前开关及本责任资格，并记录所见配置版本。发布控制以后，管理界面展示各工作者已确认版本、最后心跳、未确认实例和当时在途数。与控制面失联的实例列入未确认集合，操作者根据最后确认版本与在途明细决定是否封闭其出口。紧急暂停还可以在出口、代理或平台端限制后续请求，已经进入这些路径的请求则列入在途明细，由恢复者按原 C 查询和核对。

停止新发送与取消业务有不同的数据迁移。暂停后，未发送责任保留待办并等待策略恢复；业务撤销则由权威核验当前状态、库存决定和外部未知，提交明确的撤销或补偿责任。运维清空队列若直接删除 T，会留下无法解释的平台效果与库存占用。可执行管理接口提供暂停、重新安排到期时间、转人工和提交证据等具体动作，每个动作保留审计轨迹。

管理面与业务面采用分别授权的网络和身份。只读诊断账号没有写入权限，恢复命令按分区、商户范围和动作授权。高影响操作可要求独立复核，并把证据和批准绑定到具体命令内容。审计保存命令前后版本、接受或拒绝原因、所引用的 C/TX、输入摘要和最终提交位置。操作失败或结果未知时，操作者以原命令 ID 查询结果。

\subsection{迁移观察与人工对账}

迁移视图按源冻结、图像生成、目标安装、归属提交、目标激活和源拒绝旧写入展示。每一阶段保存迁移 ID、源目标组、源目标 epoch、图像摘要、应用位置、未决 TX 与责任数量。卡在 prepared 时，恢复者先查原协调组决定；卡在目标校验时，核对 manifest、块摘要与图像内容；归属已提交而请求还到旧位置时，入口刷新路由。每个动作的接受结果写回原迁移记录，接管者从该记录继续。

人工对账通过权威写入接口进入同一观察逻辑。操作人员提交原 C、平台对象编号、金额、币种、来源文件或查询凭证及原因。证据读取器验证来源和关联，权威检查原输入、已有结果与当前资格，再保存观察。原 C 已确定为 X1 时，另一份声称 X2 的结果产生冲突。调查者读取两份原始证据，核对商户范围、C 关联和平台明细，分清错配对象、重复导入或实际多笔收款，再按对应业务责任处理。

\CTwentyOneFigure{ch21-20}{平台账单、权威记录和履约事实经固定来源对齐，差异进入受控的恢复命令}{fig:ch21-reconcile}

批量对账先固定来源截止和版本。例如读取平台 10 月 17 日已结算账单快照 A17，以及购买权威应用到位置 L9000 的图像。匹配键优先为保存的 provider ID、商户范围和原 C 关联，金额与币种继续核对。得到“平台有、权威无”“权威成功、账单暂无”“金额不同”和“一致”等差异。账单时点、结算时点与收款时点不同，差异解释考虑各自范围。

若 A17 中出现 X1，而旧权威图像没有 C1，恢复者沿平台关联、受理日志和不可回退身份来源补回原 S、C1 与输入，再把 X1 作为原责任的观察提交。若权威有 X1 而账单暂未包含，则读取账单结算范围与生成时点，确定这笔收款属于哪次后续账单。对账修复重新进入原责任，退款则创建独立 RF，并按成功金额与未知占额检查原收款上限。批量作业保存输入快照、批次身份、差异结果和修复命令关联，重跑时读取原批次进度继续处理。

\subsection{备份恢复中的历史与责任连续性}

单副本从同组多数恢复，得到的是既有提交历史。整个权威从旧备份重建，可能丢掉已经对外产生效果的后缀。假设备份止于 C1 受理之前，平台随后完成 X1，灾难后只恢复旧备份，系统内已经没有原 S 与 \code{pay-C1} 的关联。沿回退计数器分配 C1 会与旧历史混合；即使取得更高 epoch、为同一 S 分配从未出现的 C2，平台也可能按新键再次收款。恢复准入先找回旧 S 的原绑定和效果，或者持久关闭该歧义范围，随后才为真正的新意图取得不会重用的编号。

\CTwentyOneFigure{ch21-11}{历史回退后先隔离恢复，补齐身份和效果关联，再按能力开放流量}{fig:ch21-backup}

业务备份包含完整分区事实、会话关闭与拒绝依据、编号分配器、未决 C/T、TX 决定与 prepared 责任、退款占额、Outbox，以及模式与配置版本。消费者备份包含 I/B/F、dispatch 意图和外部关联。平台、消息与对象存储按各自范围保存材料，各域的恢复点可以不同。manifest 固定完整块集合、摘要和来源应用位置，恢复者校验后才把图像纳入恢复环境。

恢复环境先封闭新受理、平台新发送和外部履约出口，保留受控查询与核验。程序读取备份边界，重放可用日志，核对关闭会话、历史拒绝和编号分配状态。从跨故障域受理记录及平台关联恢复旧 S、原 C 和原结果，或保存覆盖歧义范围的关闭依据；另外核验新意图的编号来源不会重用历史编号。恢复网络分别授权只读查询与效果发送：查询角色通过原商户和环境读取收款对象，消费者按原发运身份读取 dispatch 结果，效果发送出口则等待相应恢复阶段开放。

开放采用分能力检查。先确认原 S 查询能读到已恢复结果、旧会话按原规则拒绝、编号不会与已出现历史冲突，再启用原键查询和结果回填。TX 决定、退款占额和消费 guard 一致后，分别恢复对应后台责任，最后逐步开放新受理和实际履约。每一步记录依据、负责人和时间，回放与核对发现矛盾时继续保持相应作用域隔离。

恢复预算把发现隔离、取回备份、重放日志、外部核对、建立多数与排空积压分别计时。设备份下载速度为 200 MiB/s、图像为 600 GiB，数据传输约需 3072 秒，即 51.2 分钟，随后还有校验、加载与日志回放。若 RTO 要求 20 分钟，部署可以预置已校验副本、按块并行下载，或先恢复关键分区以减少首批取回数据。恢复过程分别记录所用图像、实际读取字节、日志起止位置和外部核对进度，查询、原责任恢复与新受理各在达到相应条件时记录完成时刻。
